\documentclass[10pt,a4paper]{amsart}
\usepackage[margin=19mm]{geometry}
\usepackage{amsmath,amssymb,mathtools}
\usepackage[scr=rsfs,cal=euler]{mathalfa}
\usepackage{xfrac}
\usepackage[unicode,hidelinks,bookmarks=false]{hyperref}
\newcommand{\ie}{i.e.\ }
\newcommand{\cf}{cf.\ }
\newcommand{\resp}{respectively\ }
\newcommand{\Subsection}{\subsection}
\providecommand{\nobreakdash}{\nobreak\hbox{-}}
\setlength{\emergencystretch}{2em}
\multlinegap0pt
\usepackage{bm}
\usepackage{bbm}
\usepackage{dsfont}
\usepackage{xspace}
\usepackage{cancel}
\usepackage{mathtools}
\usepackage{amsthm}
\makeatletter
\@ifundefined{theorem}{\newtheorem{theorem}{Theorem}}{}
\@ifundefined{lemma}{\newtheorem{lemma}[theorem]{Lemma}}{}
\makeatother
\makeatletter
\newcommand{\Sym}{\mathrm{Sym}}
\expandafter\ifx\csname mylist\endcsname\relax
\newenvironment{mylist}{\begin{list}{}{
\setlength{\parskip}{0mm}
\setlength{\topsep}{2mm}
\setlength{\parsep}{0mm}
\setlength{\itemsep}{0.5mm}
\setlength{\labelwidth}{7mm}
\setlength{\labelsep}{3mm}
\setlength{\itemindent}{0mm}
\setlength{\leftmargin}{12mm}
\setlength{\listparindent}{6mm}
}}{\end{list}}
\fi
\expandafter\ifx\csname proofof\endcsname\relax
\newenvironment{proofof}[1]{\normalsize {\it Proof of #1}.}{{\hfill $\Box$}}
\fi
\makeatother
\title[Quotients of permutation groups by nonabelian minimal normal]{Quotients of permutation groups by nonabelian minimal normal subgroups}
\author{Derek Holt}
\date{Source: arXiv:2405.19940v2 (CC BY 4.0)}
\begin{document}
\maketitle

\noindent\emph{Source and licence.} Quotients of permutation groups by nonabelian minimal normal subgroups. Version arXiv:2405.19940v2 (CC BY 4.0), distributed under the Creative Commons Attribution 4.0 International licence, as recorded in the arXiv licence metadata for this exact version and shown on the abstract page https://arxiv.org/abs/2405.19940v2. Full rights notice: licenses/arxiv-2405.19940-CC-BY-4.0.txt.\par\smallskip

\noindent\textit{Editorial scope.} Counted unit: the Lemma whose source label is \texttt{None} in the article (source0.tex lines 301--302, source label \texttt{None}), delivered together with the article's own complete original proof of it (source0.tex lines 303--331, one \texttt{proof} environment of 29 lines). No mathematics is written, repaired or re-derived: statement and proof are the article's own text line for line. Required context retained uncounted: lem:ntproj (lines 157--163). Every definition, display and label of those lines is retained unchanged. Omitted from this package and not redistributed: 1--156, 164--300, 332--512. Nothing inside a retained excerpt is omitted or altered: every retained body is delivered exactly as the article's lines are written, so no omission receipt of any kind accompanies this package. Citations: the retained excerpts carry no citation command at all, so no bibliography entry of the article is transcribed and no reference is redistributed. Reference adaptation: none was needed and none was made; every reference inside the delivered unit resolves inside it, so no unresolved reference or citation is produced. Printed slips kept literal: no printed word, symbol, hypothesis or number of the counted statement or of its proof was altered. Layout and class: the article's own class is \texttt{article} with packages including amsmath, amsfonts, amsthm, amssymb, latexsym, xspace, xcolor, enumitem, graphicx; the wrapper is the lane's pinned \texttt{amsart} editorial wrapper at 10pt on A4, which restates the theorem-like environments the retained text needs and adds the article's own macro definitions that the retained text needs (\texttt{\textbackslash Sym}), with extra wrapper packages (bm, bbm, dsfont, xspace, cancel, mathtools). The delivered numbering is the unit's own. Assets: no retained span contains a figure environment, a graphics-inclusion command, a PDF-page inclusion, a table or a TikZ picture; no image file is copied into this package, the package declares no asset dependency and no image file is redistributed with it. Licence: version arXiv:2405.19940v2 of this article is distributed under the Creative Commons Attribution 4.0 International licence, as recorded in the arXiv licence metadata for this exact version and shown on the abstract page https://arxiv.org/abs/2405.19940v2. A full rights notice is delivered at licenses/arxiv-2405.19940-CC-BY-4.0.txt. Characters: every retained excerpt is pure ASCII, so the delivered engine, XeLaTeX with its default Latin Modern text font, reports no missing character.
\begin{lemma}\label{lem:ntproj}
Let $N \unlhd G$ with $N$ a direct product of nonabelian
simple groups,
and let $K$ be a normal subgroup of $G$ with $K \cap N = 1$. Let $L$ be a normal
subgroup of $G$ with $L \le KN \cong K \times N$. Then $L$ contains all simple
direct factors of $N$ onto which elements of $L$ project non-trivially.
\end{lemma}

\begin{lemma} $G$ is transitive on $\Omega$.
\end{lemma}

\begin{proof}
Since $N$ is a minimal normal subgroup of $G$, it acts either
trivially or faithfully on each orbit of $G$.
Suppose that $G$ is intransitive on $\Omega$, and let $\Delta$ be an
orbit of $G$ on $\Omega$ on which $N$ acts nontrivially, and hence faithfully.
Let $\Gamma = \Omega \setminus \Delta$.  By induction there are homomorphisms
$\rho_\Delta: G^\Delta \to \Sym(\Delta')$ with $\ker \rho_\Delta = N^\Delta$ and
$\rho_\Gamma: G^\Gamma \to \Sym(\Gamma)$ with $\ker \rho_\Gamma = N^\Gamma$,
where $\Delta'$ is a set with $|\Delta'|<|\Delta|$.
Define $\rho:G \to \Sym(\Delta' \cup \Gamma)$ by
$\rho(g) := \rho_1(g^\Delta)\rho_2(g^\Gamma)$.
We claim that $\ker \rho = N$, which contradicts the assumption that $G$ with
normal subgroup $N$ is a counterexample to the theorem.

Clearly $N \le \ker \rho$ so suppose that $g \in \ker \rho$.
Then $g^{\Delta} \in \ker \rho_1 = N^\Delta$, so there exists
$n_1 \in N$ with $g^\Delta = n_1^\Delta$, and then $gn_1^{-1}$
is in the kernel $K_1$ of the action of $G$ on $\Delta$.
Since $(gn_1^{-1})^\Gamma \in \ker \rho_2 = N^\Gamma$, there exists
$n_2 \in N$ with $gn_1^{-1}n_2^{-1}$ in the kernel $K_2$ of the action of
$G$ on $\Gamma$, and we also have
$gn_1^{-1}n_2^{-1} \in K_1 N = K_1 \times N$, because $N$ acts
faithfully on $\Delta$.
Now, by applying Lemma~\ref{lem:ntproj} to $K_1$, $N$, and
$L:=K_2 \cap K_1 N$, we find that any simple direct factor of $N$
onto which $n_2$ projects non-trivially lies in $K_2$, and so
$n_2 \in K_2$ and hence $gn_1^{-1} \in K_1 \cap K_2=1$,
so $g=n_1 \in N$.  This proves the claim and the lemma.
\end{proof}

\end{document}
